\begin{helper}
For a fixed split-blocker witness surface, endpoint identification is a
separate obligation from adjacent monotonicity.  Let \(S_j(\omega)\) and
\(T_j(\omega)\) be the complete source and target mixed-stage events, let
\(C_\omega,C'_\omega\) be the paired cells, and let \(H(j,\omega)\) be a
normalized stage-mass family.  If the mass of \(H(0,\omega)\) is the mass of
\(S_0(\omega)\cap C_\omega\), the mass of \(H(n,\omega)\) is the mass of
\(T_n(\omega)\cap C'_\omega\), and the endpoint set identities
\[
S_0(\omega)\cap C_\omega=E_{\rm src}(\omega)\cap C_\omega,\qquad
T_n(\omega)\cap C'_\omega=E_{\rm tgt}(\omega)\cap C'_\omega
\]
hold for every \(\omega\), then the normalized endpoint masses are exactly
the source and target endpoint masses.  This includes the case \(n=0\):
there is no adjacent step to invoke, only the two endpoint set identities.
\end{helper}
\begin{proof}
Fix \(\omega\).  The source mass hypothesis gives
\[
\operatorname{ofReal}(w(\omega)H(0,\omega))
  =\nu(S_0(\omega)\cap C_\omega).
\]
Substituting the source endpoint set identity
\(S_0(\omega)\cap C_\omega=E_{\rm src}(\omega)\cap C_\omega\) gives the
claimed source endpoint formula.  The target side is identical: start from
\[
\operatorname{ofReal}(w(\omega)H(n,\omega))
  =\nu'(T_n(\omega)\cap C'_\omega)
\]
and substitute
\(T_n(\omega)\cap C'_\omega=E_{\rm tgt}(\omega)\cap C'_\omega\).
No comparison between adjacent indices is used, so the argument is valid
when \(n=0\) as well as when \(n>0\).
\end{proof}
